\infographicchapter{Fluoroscopy}{fluroscopy.png}
\label{ch:infographic-fluroscopy}

\section{Overview}

Fluoroscopy is a live, moving X-ray exam---often described as a real-time X-ray movie of the inside of your body. Instead of taking a single still X-ray snapshot, the machine sends rapid, brief pulses of X-rays through your body onto a digital flat-panel detector, showing bones, organs, swallowed contrast liquid, or medical instruments moving live on a video screen.

Doctors across many specialties rely on fluoroscopy to watch body functions in action and guide tiny tools safely. Common uses include swallowing and stomach exams (such as barium swallows, upper GI series, and barium enemas), heart catheterizations and blood-vessel exams, bladder and fertility tube tests (voiding cystourethrography and hysterosalpingography), lining up broken bones during orthopedic surgery, and guiding spine injections or pain procedures.

Its greatest strength is instant, real-time video of motion and live instrument guidance. Its limits include showing less soft-tissue detail than CT or MRI and flattening three-dimensional body parts onto a single two-dimensional screen. Because fluoroscopy uses ongoing X-ray radiation, doses can be higher than a single plain X-ray---especially during long procedures---so care teams strictly follow ALARA (As Low As Reasonably Achievable) safety practices, using short pulsed beams, tight focus, protective lead shielding, and continuous dose tracking to protect both patients and staff.

\section{Underlying Physical Principles}

\begin{itemize}
    \item \textbf{Pulsed X-Ray Movies:} An X-ray tube generates a controlled beam (typically 60--120 kVp) that pulses rapidly at 1 to 30 frames per second. Automatic brightness controls adjust the beam gently as you move so the live picture stays clear while keeping radiation as low as possible.
    \item \textbf{How Contrast Liquids and Air Outline Hollow Organs:} Dense bones block more X-rays than soft organs or air, creating natural grayscale contrast. Because hollow organs like the stomach, bowel, or blood vessels look faint on regular X-rays, doctors often use safe radiopaque ``positive'' contrast liquids---such as insoluble chalky barium sulfate or clear iodine contrast---that strongly absorb X-rays and outline those passages in sharp detail. For digestive exams, teams also pair contrast liquid with ``negative'' contrast like air or carbon dioxide gas (double-contrast imaging) to gently expand the stomach or bowel and highlight delicate surface folds.
    \item \textbf{Digital Flat-Panel Detectors:} Modern fluoroscopy units use sleek flat digital detectors that catch transmitted X-rays and convert them directly into crisp, undistorted live video without the curved edge warping of older tube-style systems.
    \item \textbf{Smart Noise Reduction and Dose Tracking:} Video computers smooth out graininess, sharpen edges, and freeze the last picture on screen (last-image-hold) so doctors can study an image without keeping the X-ray beam running. Built-in sensors track radiation dose continuously throughout the exam.
\end{itemize}

\section{Imaging Mechanism}

\begin{enumerate}
    \item \textbf{Lining Up the Table and Detector:} You stand or lie on a motorized X-ray table (or beneath a C-shaped mobile scanner) with the X-ray tube usually positioned beneath the table and the flat detector brought close to your body. Motorized lead shutters narrow the beam tightly to the exact area being checked.
    \item \textbf{Activating Live Pulsed X-Rays:} When the doctor presses a foot pedal, the machine fires brief X-ray pulses while automatically fine-tuning brightness as contrast liquid flows or as the table gently tilts.
    \item \textbf{Capturing Digital Video:} X-rays passing through your body strike the flat digital detector, which turns the pattern into high-detail digital video frames many times each second.
    \item \textbf{Real-Time Image Cleanup:} Built-in processors balance brightness, reduce graininess, sharpen fine edges, and save key still snapshots (spot images) whenever the doctor marks an important moment.
    \item \textbf{Live Display and Dose Readout:} Smooth live video appears immediately on bright room monitors, freezing the last frame whenever the pedal is released while displaying live beam time and radiation dose totals.
\end{enumerate}

\section{Key Technical Parameters}

\noindent\begin{tabularx}{\textwidth}{@{}l X@{}}
\toprule
\textbf{Parameter} & \textbf{Typical Value / Description} \\
\midrule
Physical Principle & Live moving X-ray video captured by a digital flat-panel detector \\
Energy / Signal Source & Pulsed or continuous X-ray tube (fixed fluoroscopy table or mobile C-arm scanner) \\
Typical Frequency / Energy & 60--120 kVp X-ray beam pulsed at 1--30 frames per second \\
Spatial Resolution & Fine detail down to roughly 0.15--0.35 mm pixel size on modern flat-panel detectors \\
Typical Scan Duration & 15--45 minutes for swallowing, GI, or urinary exams; 45--90 minutes for complex guided procedures \\
Contrast Agent & Swallowed or rectal barium liquid, water-soluble iodine liquid, or IV iodine contrast depending on the exam \\
Ionizing Radiation & Yes (typically 1--6 mSv for GI studies; higher localized skin doses during long procedures; strictly tracked under ALARA) \\
\bottomrule
\end{tabularx}

\section{Common Symptoms That Prompt This Scan}
\begin{itemize}
    \item \textbf{Trouble Swallowing or Choking on Food:} Feeling like food sticks in the throat or chest, coughing during meals, or persistent acid reflux and heartburn.
    \item \textbf{Ongoing Stomach or Bowel Symptoms:} Unexplained nausea, vomiting, chronic diarrhea, constipation, or abdominal pain needing a live look at how the digestive tract moves.
    \item \textbf{Joint Pain Needing Guided Injection:} Stubborn hip, shoulder, or spine pain where a doctor needs live X-ray guidance to place numbing or steroid medicine accurately.
    \item \textbf{Fertility or Urinary Tract Concerns:} Difficulty conceiving (checking fallopian tube openness) or recurrent urinary infections and leakage (checking bladder emptying).
\end{itemize}

\section{Clinical Usages}

\subsection{Primary Usages}

\begin{itemize}
    \item \textbf{Swallowing and Digestive Tract Exams:} Watches swallowing mechanics live (modified barium swallow), checks esophagus movement, examines the stomach and small bowel, and looks for leaks after digestive surgery.
    \item \textbf{Heart and Blood-Vessel Procedures:} Guides thin tubes (catheters) live during coronary angiograms, heart stent placements, rhythm procedures, pacemaker lead placements, and leg-artery treatments.
    \item \textbf{Needle and Pain-Management Guidance:} Helps doctors guide needles precisely for spinal taps (lumbar puncture), spine dye exams (myelography), epidural steroid injections, joint blocks, and central IV line placements.
    \item \textbf{Orthopedic Surgery Support:} Mobile C-arm units show live bone alignment in the operating room while surgeons set fractures, place rods, screws, or plates, and align joint replacements.
    \item \textbf{Bladder and Fertility Exams:} Checks for backward urine flow toward the kidneys in children or adults (voiding cystourethrography, or VCUG), maps urinary tubes, and checks whether fallopian tubes are open (hysterosalpingography, or HSG).
\end{itemize}

\subsection{Secondary Usages}

\begin{itemize}
    \item \textbf{Checking Breathing Muscle Motion:} A quick ``sniff test'' watches your diaphragm move as you sniff sharply to check for nerve weakness or paralysis on one or both sides.
    \item \textbf{Checking Medical Tubes and Devices:} Confirms that feeding tubes, brain shunts, kidney stents, or gastric bands sit in the right place and flow freely.
    \item \textbf{Joint Injections and Finding Foreign Objects:} Guides contrast dye into joints before an MRI or CT arthrogram and helps surgeons locate swallowed or embedded metallic objects.
\end{itemize}

\section{Advantages}

\begin{itemize}
    \item \textbf{Watches Body Functions Live:} Captures real-time motion---such as swallowing, stomach emptying, diaphragm movement, and contrast flow---that still pictures cannot show.
    \item \textbf{Precise Live Tool Guidance:} Lets doctors watch guidewires, catheters, stents, needles, and orthopedic screws move millimeter by millimeter for safer procedures.
    \item \textbf{Flexible and Widely Available:} Available both in dedicated X-ray suites and on rolling C-arm machines that move easily into operating rooms and procedure suites.
    \item \textbf{Seamless Shift from Finding to Treating:} Lets doctors diagnose a blockage or leak and treat it right away in the same session.
\end{itemize}

\section{Disadvantages}

\begin{itemize}
    \item \textbf{Ongoing X-Ray Radiation:} Long, complex procedures can deliver noticeable radiation doses to patients and staff, carrying small risks of temporary skin redness or hair loss at high doses, long-term cancer risk, or eye-lens clouding if safety rules are not followed.
    \item \textbf{Limited Soft-Tissue Detail:} Soft organs look faint compared with CT, MRI, or ultrasound, so contrast liquids are often needed to outline hollow passages.
    \item \textbf{Flat Two-Dimensional View:} Flattening 3D anatomy onto a flat screen causes overlapping bones or bowel loops and can make curved vessels look shorter than they really are.
    \item \textbf{Contrast Side Effects:} Iodine contrast carries small risks of allergic reactions or kidney stress, and barium must never be used if a bowel tear is suspected because leaking barium can irritate the belly cavity.
    \item \textbf{Pregnancy Precautions:} Elective fluoroscopy is generally avoided during pregnancy to protect the developing baby, and non-X-ray options like ultrasound or MRI are preferred whenever possible.
\end{itemize}

\section{Contraindications}

\subsection{Absolute}

\begin{itemize}
    \item \textbf{Barium and Severe Contrast Risks:} Swallowed or rectal barium must never be used if a hole or leak in the digestive tract---or a high risk of inhaling liquid into the lungs---is suspected (safe water-soluble contrast is used instead). A past severe allergic reaction without premedication rules out elective IV iodine contrast.
\end{itemize}

\subsection{Relative}

\begin{itemize}
    \item \textbf{Pregnancy:} Because fluoroscopy uses X-rays, elective belly or pelvic fluoroscopy is generally postponed during pregnancy unless urgently needed, with ultrasound or MRI used instead whenever possible.
    \item \textbf{Procedure-Specific Concerns:} Active pelvic infection (for fertility HSG exams), severe uncontrolled bleeding disorders (for invasive needle or catheter procedures), advanced kidney disease (for IV iodine contrast), or inability to swallow safely or hold required positions.
\end{itemize}

\section{Patient Safety Guidelines}

\begin{itemize}
    \item \textbf{Keeping Radiation as Low as Possible (ALARA):} Teams use low-rate pulsed X-rays, freeze the last image on screen instead of running live beams, narrow the beam tightly to the target area, keep beam-on time short, and position the detector close to your body while keeping the X-ray tube farther away.
    \item \textbf{Tracking Skin Dose During Long Procedures:} Modern units continuously display both cumulative air kerma (which estimates concentrated skin exposure) and dose-area product (DAP, which tracks total radiation delivered across the beam area). During long interventions, doctors vary the camera angle so radiation is not focused on one patch of skin, and schedule a skin checkup if high dose thresholds (such as 5 Gy) are reached.
    \item \textbf{Protecting Pregnancy:} All patients who could be pregnant are screened before X-ray exposure. Non-urgent exams are postponed or replaced with ultrasound or MRI, and a radiation physicist helps minimize baby exposure if an urgent exam is required.
    \item \textbf{Safe Contrast Choices:} Staff screen for past contrast allergies and kidney disease before IV iodine dye, and automatically choose gentle water-soluble contrast instead of barium whenever a digestive leak or choking risk is suspected.
    \item \textbf{Tested Equipment and Smart Shielding:} X-ray units undergo regular safety checks and dose calibration, staff wear protective aprons and glasses, and lead shields are kept out of the main X-ray beam path so the machine does not accidentally boost radiation output.
\end{itemize}

\section{Prerequisite Tests}

\begin{itemize}
    \item \textbf{No Routine Lab Tests for Simple Exams:} Standard swallowing or stomach fluoroscopy studies need no blood work, though reviewing past X-rays or scans helps your radiologist choose the safest contrast liquid.
    \item \textbf{Kidney Blood Test (for IV Contrast):} If IV iodine contrast is planned and you have kidney risk factors (such as kidney disease, diabetes, or older age), a recent creatinine/eGFR blood test is checked first.
    \item \textbf{Clotting Blood Tests (for Invasive Procedures):} If fluoroscopy is guiding a needle or catheter procedure with bleeding risk (like a spinal tap, myelogram, or vascular catheterization), clotting numbers (INR/PT) and platelets may be checked.
    \item \textbf{Pregnancy Check:} Patients who could be pregnant are screened by history---and a quick urine or blood pregnancy test when needed---before abdominal, pelvic, or prolonged fluoroscopy.
\end{itemize}

\section{Preparing for the Scan}

\begin{itemize}
    \item \textbf{Fasting or Bowel Prep (Depending on the Exam):} For an upper GI or esophagus study, do not eat or drink for 6--8 hours beforehand (usually after midnight). For a barium enema, follow the clear-liquid diet and bowel-cleansing prep provided by your clinic. Non-digestive studies (like joint injections) usually need no fasting unless sedation is planned.
    \item \textbf{Comfortable Clothes and Removing Metal:} Wear loose, comfortable clothing and expect to change into a hospital gown. Remove jewelry, glasses, belts, dentures (for swallowing studies), and any metal objects that could block the X-ray view.
    \item \textbf{Medicines, Allergies, and Past X-Rays:} Bring a full list of your medications, and do not stop taking any prescribed medicines unless your provider tells you to (ask ahead if you need to pause blood thinners before a needle or catheter procedure). Tell the team about any contrast, iodine, or shellfish allergies, asthma, kidney disease, suspected bowel leak, and all past X-ray exams so they can factor in your cumulative radiation history.
    \item \textbf{Pregnancy Notice and Cycle Timing:} Always tell the technologist if you might be pregnant. For a fertility HSG test, schedule the exam right after your period stops and before ovulation (typically days 7--10 of your cycle). Plan to drink extra water afterward to help clear barium smoothly.
\end{itemize}

\section{During Scan}

\begin{itemize}
    \item \textbf{Moving Table and Positioning:} You will stand or lie on a motorized X-ray table while the camera tower or C-arm moves smoothly around you. Stay still during live capture and roll gently onto your side or stomach when guided by the team.
    \item \textbf{Drinking, Breathing, and Moving on Cue:} Follow simple spoken cues---such as taking sips of chalky barium or clear contrast liquid during swallowing or stomach exams, holding your breath briefly for crisp still shots, or gently bearing down when asked.
    \item \textbf{What You Will Feel and Setting:} Exams take place in an outpatient radiology center or during a hospital stay, and the X-rays themselves are completely painless. You will hear soft mechanical hums as the table tilts or the detector glides, and you may feel mild belly fullness from swallowed or rectal contrast. For catheter or needle procedures (often placed through a site such as the groin or elbow), you may feel a brief pinch from numbing medicine or receive relaxing pain medication---while surgical or stent-placement procedures may use general anesthesia.
    \item \textbf{Easy Communication:} The radiologist and technologist stay right there talking with you throughout the exam. Speak up right away if you feel dizzy, nauseated, short of breath, sharp pain, or coughing while swallowing contrast.
    \item \textbf{How Long It Takes:} Most diagnostic swallowing, stomach, or bladder fluoroscopy exams take 15--45 minutes. Complex catheter or pain-management procedures may take 45--90 minutes.
\end{itemize}

\section{Normal Appearance}

\begin{itemize}
    \item \textbf{Smooth Swallowing and Esophagus:} Liquid moves smoothly from the mouth down the throat with no spillover into the windpipe, the esophagus walls look smooth and even, and steady wave-like squeezes carry contrast cleanly into the stomach.
    \item \textbf{Healthy Stomach and Bowel:} Contrast coats the stomach and intestines evenly with no ulcers, bumps, or sudden pinch points, emptying smoothly through neat natural bowel folds.
    \item \textbf{Normal Diaphragm, Bladder, and Fertility Tubes:} Both sides of the breathing muscle (diaphragm) move evenly downward on a sharp sniff; the bladder looks smooth with no urine backing up toward the kidneys on a VCUG; and the uterus looks smoothly triangular with dye flowing freely out both fallopian tubes on an HSG.
    \item \textbf{Aligned Bones and Devices:} Broken bones sit cleanly lined up, and catheters, stents, or feeding tubes sit accurately in place with no leaking contrast.
\end{itemize}

\section{Observing for Abnormal Patterns}

\begin{itemize}
    \item \textbf{Narrowings and Lumps:} A smooth, tapered ``bird-beak'' narrowing at the bottom of the esophagus points to achalasia (a tight swallowing valve), while a jagged, abrupt ``apple-core'' narrowing on a barium swallow or enema is a classic sign of a tumor. Dark round spots inside bright contrast can signal polyps, masses, or stones.
    \item \textbf{Contrast Leaks and Abnormal Tunnels:} Bright contrast escaping outside the stomach, bowel, urinary tract, or a blood vessel confirms a tear, surgical leak, or active bleeding. Thin contrast tracks branching to nearby organs or skin outline abnormal tunnels (fistulas).
    \item \textbf{Swallowing and Movement Problems:} Contrast slipping into the airway on a modified barium swallow confirms swallowing difficulty and aspiration risk. Uncoordinated, twisting esophagus contractions (``corkscrew esophagus'') signal esophageal spasm, and contrast backing up from the bladder into the kidney tubes grades urinary reflux.
    \item \textbf{Paralyzed Breathing Muscle:} On a sniff test, if one side of the diaphragm moves upward while the healthy side moves downward, it confirms paralysis of that side's breathing nerve.
    \item \textbf{Harmless Look-Alikes:} Overlapping loops of bowel or bones can look like narrowings or lumps until the doctor turns you to a different angle or presses gently to separate them. Tiny air bubbles or leftover stool specks can also look like polyps on double-contrast studies.
\end{itemize}

\section{Dealing with Uncertainties}

\begin{itemize}
    \item \textbf{Overlapping Shadows:} Because fluoroscopy flattens 3D body parts onto a 2D screen, overlapping bowel loops or angled blood vessels can look pinched or swollen. Turning to side angles, gently pressing the area, and snapping still spot images help separate true problems from overlap.
    \item \textbf{Bubbles and Fast Flow:} Leftover stool, mucus, or air bubbles can look like tiny polyps, while fast liquid flow or awkward positioning can briefly miss intermittent reflux, subtle swallowing problems, or tiny surgical leaks.
    \item \textbf{Tell-Tale Narrowings That Overlap:} Smooth versus jagged narrowing alone cannot always tell a harmless acid or scar narrowing apart from a tumor or an outside organ pressing on the bowel.
    \item \textbf{Confirming with Other Tests:} Unclear fluoroscopy findings are commonly followed up with a direct camera scope (endoscopy) and biopsy, a 3D CT or MRI scan for deeper organs, or specialized swallowing pressure tests (manometry) and speech-therapy review.
\end{itemize}

\section{Augmenting Methods}

\begin{itemize}
    \item \textbf{Endoscopy and Biopsy:} Lets doctors look directly at any narrowing, ulcer, polyp, or mass seen on a barium study and take tissue samples for an exact answer.
    \item \textbf{CT and MRI Scans:} Add 3D cross-sectional detail of organ walls, deep fluid pockets, and blood vessels that flat fluoroscopy cannot show.
    \item \textbf{Swallowing Pressure Tests and Camera Swallowing Exams:} Esophageal pressure testing (manometry) measures exact muscle squeezes for achalasia or spasm, while a tiny nose camera (FEES) complements modified barium swallows for swallowing care.
    \item \textbf{3D Cone-Beam CT and Ultrasound Guidance:} Combining 3D rotational X-ray scans or live ultrasound during fluoroscopy procedures adds soft-tissue clarity and helps keep radiation doses lower.
\end{itemize}

\section{Comparison with Alternative Modalities}

\begin{itemize}
    \item \textbf{Compared with CT Scans:} CT gives far sharper 3D cross-sectional views of soft tissues without overlapping shadows, while fluoroscopy is better for watching live motion like swallowing, bowel movement, and real-time catheter or screw placement.
    \item \textbf{Compared with Endoscopy:} Endoscopy provides direct full-color views of the inner lining and immediate biopsies, while fluoroscopy checks swallowing and muscle motion noninvasively, maps tight narrowings a scope cannot pass, and spots surgical leaks safely.
    \item \textbf{Compared with Ultrasound:} Ultrasound uses zero radiation and excels at guiding needles into soft tissues near the surface or belly, while fluoroscopy sees right through bone and bowel gas and gives a broad view for blood vessels, bones, and digestive transit.
    \item \textbf{Compared with Standard Still X-Rays:} Plain X-rays use less radiation for a single frozen snapshot, while fluoroscopy adds continuous live motion for functional exams and hands-on procedure guidance.
\end{itemize}

\section{Follow-up Steps}

\begin{itemize}
    \item \textbf{Drinking Water and Clearing Contrast:} After a barium digestive study, drink plenty of water for a day or two to prevent constipation; pale or white stools are completely normal for 1--3 days. After IV iodine contrast or an invasive procedure, staff monitor your puncture site and check kidney function if you are at higher risk.
    \item \textbf{Radiation Dose Record and Skin Check:} Your total radiation dose is recorded in your chart. If a long, complex procedure exceeds high dose thresholds (such as 5 Gy), your team will let you know and schedule a quick skin check in 2--4 weeks to watch for mild redness or irritation.
    \item \textbf{Next Steps for Findings:} Suspicious narrowings or bumps are followed up with endoscopy and biopsy or a CT/MRI scan; after a modified barium swallow, a speech-language pathologist uses the video to recommend safer foods, liquids, or swallowing exercises.
    \item \textbf{Prompt Results:} Urgent findings---such as a bowel leak or severe choking risk---are shared with your doctor immediately, and repeat fluoroscopy is ordered only when truly needed to confirm healing or device position.
\end{itemize}

\section{Key References}

\begin{itemize}
    \item American College of Radiology (ACR). \textit{ACR--AAPM Technical Standard for Management of the Use of Radiation in Fluoroscopic Procedures}. Reston, VA: ACR, 2023.
    \item National Council on Radiation Protection and Measurements (NCRP). \textit{NCRP Report No. 168: Radiation Dose Management for Fluoroscopically-Guided Interventional Medical Procedures}. Bethesda, MD: NCRP, 2010.
    \item American College of Radiology (ACR). \textit{ACR Manual on Contrast Media}. Reston, VA: ACR, 2024.
    \item Radiological Society of North America (RSNA) and ACR. \textit{RadiologyInfo Clinical Patient and Safety Resources: Fluoroscopy}. 2024.
    \item National Library of Medicine (MedlinePlus). \textit{Fluoroscopy: MedlinePlus Medical Test}. National Institutes of Health (NIH). Available at: \url{https://medlineplus.gov/lab-tests/fluoroscopy/}.
    \item Radiopaedia.org. \textit{Fluoroscopy}. Radiopaedia. Available at: \url{https://radiopaedia.org/articles/fluoroscopy}.
\end{itemize}

\section{Further Reading}
\begin{itemize}
  \item Wikipedia contributors. \textit{Fluoroscopy}. Wikipedia, The Free Encyclopedia. Available at: \url{https://en.wikipedia.org/wiki/Fluoroscopy}.
\end{itemize}
